%%=============================================================================
%% Voorwoord
%%=============================================================================

\chapter*{\IfLanguageName{dutch}{Woord vooraf}{Preface}}%
\label{ch:voorwoord}

%% TODO:
%% Het voorwoord is het enige deel van de bachelorproef waar je vanuit je
%% eigen standpunt (``ik-vorm'') mag schrijven. Je kan hier bv. motiveren
%% waarom jij het onderwerp wil bespreken.
%% Vergeet ook niet te bedanken wie je geholpen/gesteund/... heeft

Deze bachelorproef over het automatiseren van CRA-compliance in CI/CD-pipelines is voor mij een logische keuze geweest. Vanwege de European Cyber Resilience Act, die vanaf 2027 van kracht wordt, staan softwarebedrijven voor een grote uitdaging om securitycontroles systematisch te integreren in hun ontwikkelproces. Dit onderwerp is niet alleen uiterst actueel, maar sluit ook perfect aan bij mijn passie voor cybersecurity, automatisatie en DevSecOps.

Tijdens mijn bachelorproef aan de HoGent  zag ik de kans om een praktische oplossing te ontwikkelen voor een reëel probleem: het ontbreken van geautomatiseerde security scanning in bestaande Jenkins pipelines bij Excentis. Het was een boeiend proces om literatuurstudie en technische implementatie te combineren tot een werkende Proof of Concept.

Graag wil ik hierbij mijn oprechte dank betuigen aan:

\textbf{Moreno Robyn}, copromotor bij Excentis, die mij dit relevante onderwerp aanbood en de praktische context binnen het bedrijf mogelijk maakte.

\textbf{Pieter-Jan Maenhaut}, promotor, die mij met constructieve feedback en waardevolle suggesties doorheen het schrijf- en verbeterproces heeft begeleid.

